% !TEX root = ../research_vs.tex

\section{Robustness}\label{sec:robust}

\subsection{Robustness: Research vs Data Mining}\label{sec:robust,disc}


% EXHIBIT-NUMBER: Figure IA.1
\begin{figure}[h]\caption{Excluding Data-Mined Predictors Correlated with Any Existing Research}
  \label{fig:decay-unspanned}
  We compare published predictors (solid) to data-mined predictors made with t-stats $>$ 2.0, but separate predictors by correlation and in-sample t-stats. Panel (a) uses the maximum pairwise correlation with any existing published predictor. Panel (b) uses the $R^2$ from regressing the data-mined return on 5 principal components of existing predictors computed using probabilistic PCA. 
  
  \vspace{2 ex}
  \centering
  \subfloat[Any Correlation\label{fig:abnormal_published_1}]{
      \includegraphics[width=0.48\textwidth]{\exhibitspath/Fig_DM_unspan_match_t_g_cor.pdf}
  }    
  \subfloat[PPCA\label{fig:abnormal_published_2}]{
      \includegraphics[width=0.48\textwidth]{\exhibitspath/Fig_DM_unspan_match_t_g_PCA.pdf}
  }
\end{figure}

% EXHIBIT-NUMBER: Table 3
\begin{table}[!h]
   \caption{Post-Sample Outperformance by Discipline and Journal Ranking}
   \label{tab:hetero-byjournal}

   \begin{singlespace}
   \noindent Strategies are normalized to have 100 bps per month performance in-sample. `Post-Sample' is the performance of predictor-papers in bps per month. `Versus Data Mining' is `Post-Sample' minus the performance of data-mined benchmarks. `Discipline' categorizes the journal of publication. `JF,' `JFE,' and `RFS' includes papers in the Journal of Finance, Journal of Financial Economics, and Review of Financial Studies. `AR,' `JAR,' and `JAE' includes papers in the Accounting Review, Journal of Accounting Research, and Journal of Accounting and Economics. CAPM and FF3 + momentum alphas use regressions specific to the sample periods (in-sample or post-sample). Standard errors (parentheses) are clustered by calendar month and predictor. Finance, and particularly top-ranked finance journals show some signs of outperformance relative to data mining, but the improvement is modest.
   \end{singlespace}
   \begin{centering}
   \vspace{0ex}
   \par\end{centering}
   \centering{}\setlength{\tabcolsep}{1.2ex} \small
   \begin{center}
   \begin{tabular}{lcccccc}
   \toprule
   & \multicolumn{2}{c}{Long-Short Return} & \multicolumn{2}{c}{CAPM Alpha} & \multicolumn{2}{c}{FF3 + Mom Alpha} \\
   \cmidrule(lr){2-3}\cmidrule(lr){4-5}\cmidrule(lr){6-7}
   & Post- & \multicolumn{1}{c}{Versus} & Post- & \multicolumn{1}{c}{Versus} & Post- & \multicolumn{1}{c}{Versus} \\
   & Sample & \multicolumn{1}{c}{Data Mining} & Sample & \multicolumn{1}{c}{Data Mining} & Sample & \multicolumn{1}{c}{Data Mining} \\
   \midrule
   \multicolumn{7}{l}{\textbf{Discipline}} \\
   Finance & 59 & 11 & 69 & 15 & 78 & 15 \\
   & (8) & (7) & (7) & (7) & (8) & (9) \\
   Accounting & 43 & -5 & 48 & -4 & 46 & -17 \\
   & (9) & (10) & (9) & (10) & (8) & (10) \\
   \addlinespace
   \multicolumn{7}{l}{\textbf{Journal Rank}} \\
   JF, JFE, RFS & 60 & 10 & 72 & 17 & 79 & 14 \\
   & (8) & (8) & (8) & (8) & (9) & (9) \\
   AR, JAR, JAE & 43 & -5 & 45 & -7 & 44 & -20 \\
   & (9) & (10) & (9) & (10) & (8) & (10) \\
   Other & 53 & 8 & 56 & 6 & 65 & 9 \\
   & (9) & (9) & (9) & (9) & (11) & (13) \\
   \bottomrule
\end{tabular}

   \end{center}
\end{table}


\clearpage

% EXHIBIT-NUMBER: Table B.1
\begin{table}[H]
    \caption{Individual data-mined returns post-sample and post-publication}\label{tab:MPdecayIndividualDM}
    \small

    The table shows regressions of individual data-mined strategies on indicators for the sample-end and publication dates of their matched academic predictors. As in Tables \ref{tab:MPdecaymain}--\ref{tab:MPdecaytimefe}, strategies have at least 60 original-sample months, match the academic predictor within 10\% on its original-sample t-statistic and 30\% on its mean return, and exclude pairs with signed-return correlations above 10\%. For computational reasons, we sample at most 100 strategies for each academic predictor.
    Scaled returns are normalized such that a strategy's mean return is 100 bps per month in the sample of the original matched predictor. Standard errors are clustered at the data-mined strategy and calendar month levels and are shown in parentheses.
    \vspace{2ex}

    \centering

    \setlength{\tabcolsep}{2.0ex}
    
\begingroup
\centering
\begin{tabular}{lcccc}
   \toprule
    & \multicolumn{2}{c}{Scaled returns} & \multicolumn{2}{c}{Unscaled returns} \\ 
                             & (1)           & (2)           & (3)           & (4)\\  
   \midrule 
   Post-Sample               & -44.3         & -36.7         & -21.3         & -19.6\\   
                             & (8.14)        & (4.50)        & (3.86)        & (2.33)\\   
   Post-Pub                  & -16.2         & 5.70          & -7.93         & 1.64\\   
                             & (8.32)        & (3.44)        & (4.20)        & (1.61)\\   
    \\
   Observations              & 4,779,040     & 4,779,040     & 4,779,040     & 4,779,040\\  
   R$^2$                     & 0.003         & 0.043         & 0.006         & 0.057\\  
   Within R$^2$              & 0.002         & 0.000         & 0.002         & 0.000\\  
    \\
   DM strategy fixed effects & $\checkmark$  & $\checkmark$  & $\checkmark$  & $\checkmark$\\   
   Month fixed effects       &               & $\checkmark$  &               & $\checkmark$\\   
   \bottomrule
\end{tabular}
\par\endgroup



\end{table}


% EXHIBIT-NUMBER: Table B.2
\begin{table}[H]
    \caption{Decay regressions using only annual accounting predictors}\label{tab:MPdecayacctonly}
    \small

    The table repeats Table \ref{tab:MPdecaymain} after restricting the academic predictors to those built from annual accounting data, the same data source as the data-mined strategies. Following Figure \ref{fig:pub-vs-dm-2}, Panel (b), we begin from the accounting category of \citet{ChenZimmermann2021} and then drop quarterly, analyst-based, and discrete signals, which leaves 66 of the 173 predictors. Post-sample is one in months after the end of the original sample and zero otherwise. Post-pub (Add'l) is one in months after publication and zero otherwise. Standard errors (parentheses) are clustered by predictor and calendar month.
    \vspace{2ex}

    \centering

    \setlength{\tabcolsep}{2.0ex}
    % GENERATED -- do not hand-edit (S3b_MPStyleDecayTables.R).
\begingroup
\setlength{\tabcolsep}{0.55ex}%
\centering
\begin{tabular}{lcccccc}
   \toprule
                    & \multicolumn{3}{c}{(a) LHS = Scaled Return} & \multicolumn{3}{c}{(b) LHS = Unscaled Return} \\
   \cmidrule(lr){2-4}\cmidrule(lr){5-7}
                    & Academic
                    & Data Mining
                    & Academic
                    & Academic
                    & Data Mining
                    & Academic \\
                    & Signal
                    & Benchmark
                    & Minus
                    & Signal
                    & Benchmark
                    & Minus \\
                    & Return
                    & (Uncorr)
                    & Data Mining
                    & Return
                    & (Uncorr)
                    & Data Mining \\
                    & (1) & (2) & (3) & (4) & (5) & (6) \\
   \midrule
   Post-Sample      & -20.4 & -38.4 & 17.9 & -7.14 & -19.4 & 12.2 \\
   SE               & (12.2) & (8.94) & (13.5) & (6.56) & (5.28) & (7.77) \\
   Post-Pub (Add'l) & -33.4 & -14.4 & -19.0 & -19.8 & -7.77 & -12.1 \\
   SE               & (13.5) & (9.16) & (15.3) & (7.77) & (5.23) & (8.35) \\
   \\
   Observations     & 41,016 & 41,016 & 41,016 & 34,842 & 34,842 & 34,842 \\
   Signals          & 66 & 66 & 66 & 57 & 57 & 57 \\
   R$^2$            & 0.003 & 0.015 & 0.001 & 0.006 & 0.020 & 0.002 \\
   Within R$^2$     & 0.001 & 0.013 & 0.000 & 0.001 & 0.006 & 0.000 \\
   \\
   Predictor F.E.   & Yes & Yes & Yes & Yes & Yes & Yes \\
   Time F.E.        & No & No & No & No & No & No \\
   \bottomrule
\end{tabular}
\par\endgroup

\end{table}

% EXHIBIT-NUMBER: Table B.3
\begin{table}[H]
    \caption{Decay regressions using only annual accounting predictors, controlling for time effects}\label{tab:MPdecayacctonlytimefe}
    \small

    The table repeats Table \ref{tab:MPdecaytimefe} on the annual accounting predictors of Table \ref{tab:MPdecayacctonly}, adding time fixed effects to the specification. It uses the same benchmark construction and the same common sample as Table \ref{tab:MPdecayacctonly}. Post-sample is one in months after the end of the original sample and zero otherwise. Post-pub (Add'l) is one in months after publication and zero otherwise. Standard errors (parentheses) are clustered by predictor and calendar month.
    \vspace{2ex}

    \centering

    \setlength{\tabcolsep}{2.0ex}
    % GENERATED -- do not hand-edit (S3b_MPStyleDecayTables.R).
\begingroup
\setlength{\tabcolsep}{0.55ex}%
\centering
\begin{tabular}{lcccccc}
   \toprule
                    & \multicolumn{3}{c}{(a) LHS = Scaled Return} & \multicolumn{3}{c}{(b) LHS = Unscaled Return} \\
   \cmidrule(lr){2-4}\cmidrule(lr){5-7}
                    & Academic
                    & Data Mining
                    & Academic
                    & Academic
                    & Data Mining
                    & Academic \\
                    & Signal
                    & Benchmark
                    & Minus
                    & Signal
                    & Benchmark
                    & Minus \\
                    & Return
                    & (Uncorr)
                    & Data Mining
                    & Return
                    & (Uncorr)
                    & Data Mining \\
                    & (1) & (2) & (3) & (4) & (5) & (6) \\
   \midrule
   Post-Sample      & -16.2 & -39.0 & 22.8 & -7.10 & -20.8 & 13.7 \\
   SE               & (14.7) & (5.06) & (15.3) & (7.27) & (5.07) & (9.33) \\
   Post-Pub (Add'l) & -14.9 & -0.613 & -14.3 & -7.79 & -1.05 & -6.74 \\
   SE               & (14.0) & (4.11) & (16.4) & (8.41) & (3.46) & (10.3) \\
   \\
   Observations     & 41,016 & 41,016 & 41,016 & 34,842 & 34,842 & 34,842 \\
   Signals          & 66 & 66 & 66 & 57 & 57 & 57 \\
   R$^2$            & 0.047 & 0.430 & 0.036 & 0.065 & 0.212 & 0.030 \\
   Within R$^2$     & 0.000 & 0.004 & 0.000 & 0.000 & 0.002 & 0.000 \\
   \\
   Predictor F.E.   & Yes & Yes & Yes & Yes & Yes & Yes \\
   Time F.E.        & Yes & Yes & Yes & Yes & Yes & Yes \\
   \bottomrule
\end{tabular}
\par\endgroup

\end{table}

\subsection{Alternative Portfolio Constructions}\label{sec:robust,port}

% EXHIBIT-NUMBER: Table B.4
\begingroup
\captionsetup{type=table}
\caption{Post-sample and post-publication decay across portfolio constructions}\label{tab:timefe-robustness}
\small

The table repeats the decay regressions of Tables \ref{tab:MPdecaymain} and \ref{tab:MPdecaytimefe} across twelve constructions of the published long-short portfolios. Within each panel, the first row uses predictor fixed effects only, as in Table \ref{tab:MPdecaymain}, and the second row adds time fixed effects, as in Table \ref{tab:MPdecaytimefe}. Post-sample is the coefficient on an indicator for months after the end of the original sample; additional post-publication is the incremental coefficient on an indicator for months after publication. Both are in basis points per month, with standard errors in parentheses. The first block uses the portfolios of the original papers, as assembled by \citet{mclean2016does} and \citet{ChenZimmermann2021}. The second block rebuilds the \citet{ChenZimmermann2021} portfolios under alternative sorting, weighting, and screening choices. The third block rebuilds them as \citet{jensen2022there} style terciles.

Post-sample decay is present under every construction. The additional post-publication decay is the more fragile estimate. Adding time fixed effects removes nearly all of it under the alternative \citet{ChenZimmermann2021} constructions of the second block, where the remaining estimates run from $+4.4$ to $-7.0$ bps per month. The \citet{jensen2022there} style terciles are the exception. There, additional post-publication decay falls but survives, at $-18.6$ bps per month for the equal-weighted terciles (Panel J) and $-21.8$ bps per month for the metadata-matched sample (Panel L), the latter more than two standard errors from zero. The \citet{mclean2016does} portfolios of Panel A likewise retain $-22.5$ bps per month, though only an approximate standard error is available there. Whether publication carries information beyond the passage of time therefore depends on how the published portfolios are rebuilt.
\vspace{2ex}

\begingroup
\small
\begin{longtable}{@{}lcccc@{}}
\toprule
Data and specification & Post-sample & \shortstack{Additional\\post-publication} & \shortstack{Mean in-sample\\return (bp/month)} & Signals \\
\midrule
\endfirsthead
\multicolumn{5}{c}{\tablename\ \thetable{} -- continued from previous page} \\
\toprule
Data and specification & Post-sample & \shortstack{Additional\\post-publication} & \shortstack{Mean in-sample\\return (bp/month)} & Signals \\
\midrule
\endhead
\midrule
\multicolumn{5}{r}{Continued on next page} \\
\endfoot
\bottomrule
\endlastfoot
\multicolumn{5}{@{}l}{\textbf{Exact published portfolios}} \\
\addlinespace[0.15em]
\multicolumn{5}{@{}l}{\textit{Panel A: MP published sample}} \\
Predictor FE & -25.8 (13.2) & -32.1 (14.3) & 58.2 & 97 \\
Predictor + time FE & -30.8 (13.7) & -22.5 ($\approx 21.0$) & 58.2 & 97 \\
\addlinespace[0.35em]
\multicolumn{5}{@{}l}{\textit{Panel B: CZ original portfolios}} \\
Predictor FE & -23.0 (11.0) & -27.6 (12.1) & 74.0 & 185 \\
Predictor + time FE & -13.8 (11.7) & -5.6 (9.2) & 74.0 & 185 \\
\addlinespace[0.4em]
\midrule
\multicolumn{5}{@{}l}{\textbf{Alternative CZ portfolio constructions}} \\
\addlinespace[0.15em]
\multicolumn{5}{@{}l}{\textit{Panel C: Deciles, equal weighted}} \\
Predictor FE & -27.5 (11.2) & -22.1 (11.9) & 91.4 & 142 \\
Predictor + time FE & -8.3 (10.9) & 3.9 (10.6) & 91.4 & 142 \\
\addlinespace[0.35em]
\multicolumn{5}{@{}l}{\textit{Panel D: Deciles, value weighted}} \\
Predictor FE & -58.6 (13.1) & -10.4 (17.4) & 92.6 & 76 \\
Predictor + time FE & -53.4 (9.7) & 2.7 (16.7) & 92.6 & 76 \\
\addlinespace[0.35em]
\multicolumn{5}{@{}l}{\textit{Panel E: Quintiles, equal weighted}} \\
Predictor FE & -25.0 (12.8) & -33.7 (13.7) & 70.2 & 149 \\
Predictor + time FE & -12.0 (13.1) & -6.0 (11.0) & 70.2 & 149 \\
\addlinespace[0.35em]
\multicolumn{5}{@{}l}{\textit{Panel F: Quintiles, value weighted}} \\
Predictor FE & -44.2 (18.6) & -31.1 (21.9) & 63.2 & 82 \\
Predictor + time FE & -36.0 (16.9) & -7.0 (16.6) & 63.2 & 82 \\
\addlinespace[0.35em]
\multicolumn{5}{@{}l}{\textit{Panel G: Exclude below NYSE 20th pct. size}} \\
Predictor FE & -40.0 (11.5) & -14.4 (14.4) & 66.8 & 129 \\
Predictor + time FE & -23.5 (11.2) & 4.4 (10.5) & 66.8 & 129 \\
\addlinespace[0.35em]
\multicolumn{5}{@{}l}{\textit{Panel H: NYSE stocks only}} \\
Predictor FE & -42.2 (11.0) & -10.5 (12.1) & 65.3 & 126 \\
Predictor + time FE & -29.7 (11.9) & 1.0 (10.5) & 65.3 & 126 \\
\addlinespace[0.35em]
\multicolumn{5}{@{}l}{\textit{Panel I: Exclude price below \$5}} \\
Predictor FE & -46.6 (9.6) & -15.3 (11.2) & 67.0 & 159 \\
Predictor + time FE & -35.1 (10.1) & 2.5 (9.4) & 67.0 & 159 \\
\addlinespace[0.4em]
\midrule
\multicolumn{5}{@{}l}{\textbf{Rebuilt as JKP-style terciles}} \\
\addlinespace[0.15em]
\multicolumn{5}{@{}l}{\textit{Panel J: JKP equal-weighted terciles}} \\
Predictor FE & -11.9 (10.5) & -32.0 (12.7) & 60.3 & 97 \\
Predictor + time FE & 1.4 (11.3) & -18.6 (10.9) & 60.3 & 97 \\
\addlinespace[0.35em]
\multicolumn{5}{@{}l}{\textit{Panel K: CZ reconstructed as equal-weighted terciles}} \\
Predictor FE & -21.4 (9.6) & -32.4 (10.5) & 61.4 & 135 \\
Predictor + time FE & -13.0 (8.4) & -8.8 (9.3) & 61.4 & 135 \\
\addlinespace[0.35em]
\multicolumn{5}{@{}l}{\textit{Panel L: CZ terciles, JKP-metadata-matched signals}} \\
Predictor FE & -13.7 (11.1) & -44.1 (12.0) & 58.0 & 69 \\
Predictor + time FE & -3.9 (10.1) & -21.8 (9.9) & 58.0 & 69 \\
\end{longtable}
{\footnotesize\emph{Note:} Each panel reports two rows, predictor fixed effects and predictor plus time fixed effects. Columns give the post-sample change and the additional post-publication change in mean returns (standard errors in parentheses), and the mean in-sample long-short return. All rows except MP are scaled by each signal's in-sample mean return; the MP rows are scaled by the grand mean across signals. $\approx$ marks an approximate standard error.\par}
\endgroup

% The generated fragment is a bare longtable with no \caption of its own, so it
% still steps the table counter. The caption above already claimed this table's
% number; give the stray increment back so the next float numbers correctly.
\addtocounter{table}{-1}
\endgroup
